\documentclass[12pt]{article}
\usepackage{amsmath}
\usepackage{amsfonts}
\usepackage{geometry}
\usepackage{multicol}
\geometry{margin=1in}

\title{Calculus 1 Survival Sheet}
\date{}
\begin{document}

\maketitle

%%%%%%%%%%%%%%%%%%%%%%%%%%%%%%%%%%%%%%%%%%%%%%%%%%%%%%
% ORIGINAL CONTENT (UNTOUCHED)
%%%%%%%%%%%%%%%%%%%%%%%%%%%%%%%%%%%%%%%%%%%%%%%%%%%%%%

\section*{Definition of the Derivative}

The derivative of a function $f(x)$ at a point $x$ is defined as:
\[
f'(x) = \lim_{h \to 0} \frac{f(x + h) - f(x)}{h}
\]

This limit, if it exists, gives the slope of the tangent line to the curve at the point $x$.

\section*{Basic Derivatives}

\begin{align*}
\frac{d}{dx}(c) &= 0 \quad \text{(constant rule)} \\
\frac{d}{dx}(x) &= 1 \\
\frac{d}{dx}(x^n) &= nx^{n-1} \\
\frac{d}{dx}(e^x) &= e^x \\
\frac{d}{dx}(\ln x) &= \frac{1}{x} \\
\frac{d}{dx}(\sin x) &= \cos x \\
\frac{d}{dx}(\cos x) &= -\sin x \\
\frac{d}{dx}(\tan x) &= \sec^2 x \\
\frac{d}{dx}(\sec x) &= \sec x \tan x \\
\frac{d}{dx}(\csc x) &= -\csc x \cot x \\
\frac{d}{dx}(\cot x) &= -\csc^2 x
\end{align*}

\subsection*{Constant Multiple Rule}
\[
\frac{d}{dx}\left[c \cdot f(x)\right] = c \cdot f'(x)
\]

\section*{Rules of Differentiation}

\subsection*{Sum and Difference Rule}
\[
\frac{d}{dx}[f(x) \pm g(x)] = f'(x) \pm g'(x)
\]

\subsection*{Product Rule}
\[
\frac{d}{dx}[f(x)g(x)] = f'(x)g(x) + f(x)g'(x)
\]

\subsection*{Quotient Rule}
\[
\frac{d}{dx}\left[\frac{f(x)}{g(x)}\right] = \frac{f'(x)g(x) - f(x)g'(x)}{[g(x)]^2}
\]

\section*{Chain Rule}
If $y = f(g(x))$, then:
\[
\frac{dy}{dx} = f'(g(x)) \cdot g'(x)
\]

\textbf{Example:} If $y = \sin(x^2)$, then:
\[
\frac{dy}{dx} = \cos(x^2) \cdot 2x
\]

\section*{Implicit Differentiation}

When $y$ is defined implicitly as a function of $x$, differentiate both sides with respect to $x$, applying the chain rule to $y$ terms.

\textbf{Example:} For the equation $x^2 + y^2 = 25$:

\begin{align*}
\frac{d}{dx}(x^2) + \frac{d}{dx}(y^2) &= \frac{d}{dx}(25) \\
2x + 2y\frac{dy}{dx} &= 0 \\
\frac{dy}{dx} &= -\frac{x}{y}
\end{align*}

\section*{Useful Derivatives}

\begin{multicols}{2}
\begin{itemize}
    \item $\frac{d}{dx}(a^x) = a^x \ln a$
    \item $\frac{d}{dx}(\arcsin x) = \frac{1}{\sqrt{1 - x^2}}$
    \item $\frac{d}{dx}(\arccos x) = \frac{-1}{\sqrt{1 - x^2}}$
    \item $\frac{d}{dx}(\arctan x) = \frac{1}{1 + x^2}$
\end{itemize}
\end{multicols}

%%%%%%%%%%%%%%%%%%%%%%%%%%%%%%%%%%%%%%%%%%%%%%%%%%%%%%
% NEW SECTIONS ADDED BELOW
%%%%%%%%%%%%%%%%%%%%%%%%%%%%%%%%%%%%%%%%%%%%%%%%%%%%%%

\section*{More Trig Derivatives}
\[
\frac{d}{dx}(arccot x)= -\frac{1}{1+x^2}, \quad
\frac{d}{dx}(arcsec x)= \frac{1}{|x|\sqrt{x^2-1}}
\]
\[
\frac{d}{dx}(arccsc x)= -\frac{1}{|x|\sqrt{x^2-1}}
\]

\section*{Basic Integration Rules}

\begin{align*}
\int x^n\,dx &= \frac{x^{n+1}}{n+1} + C \quad (n\neq -1)\\
\int \frac{1}{x} \, dx &= \ln|x| + C \\
\int e^x \, dx &= e^x + C \\
\int a^x \, dx &= \frac{a^x}{\ln a} + C
\end{align*}

\subsection*{Trig Integrals}
\begin{align*}
\int \sin x\,dx &= -\cos x + C \\
\int \cos x\,dx &= \sin x + C \\
\int \sec^2 x\,dx &= \tan x + C \\
\int \csc^2 x\,dx &= -\cot x + C \\
\int \sec x \tan x\,dx &= \sec x + C \\
\int \csc x \cot x\,dx &= -\csc x + C
\end{align*}

\section*{Integration by $u$-Substitution}

If $u=g(x)$, then
\[
\int f(g(x))g'(x)\,dx = \int f(u)\,du.
\]

Steps:
\begin{enumerate}
\item Choose $u$ (inside part).
\item Compute $du$.
\item Rewrite integral in $u$.
\item Integrate.
\item Substitute back.
\end{enumerate}

\section*{Definition of $\ln x$ and $e^x$}
\[
\ln x = \int_1^x \frac{1}{t}\,dt.
\]
The function $e^x$ is the inverse of $\ln x$.

\section*{Summation Rules}
\[
\sum_{i=1}^{n} i = \frac{n(n+1)}{2}
\]

\[
\sum_{i=1}^{n} i^{2} = \frac{n(n+1)(2n+1)}{6}
\]

\[
\sum_{i=1}^{n} i^{3} = \left( \frac{n(n+1)}{2} \right)^{2}
\]

% Linearity
\[
\sum_{i=1}^{n} (a f(i) + b g(i)) 
= a\sum_{i=1}^{n} f(i) + b\sum_{i=1}^{n} g(i)
\]

% Constant Rule
\[
\sum_{i=1}^{n} c = cn
\]

% Shifting the Index
\[
\sum_{i=k}^{n} f(i) 
= \sum_{j=1}^{n-k+1} f(j+k-1)
\]
\section*{Riemann Sums}
\[
\text{Left sum: } L_n = \sum_{i=1}^n f(x_{i-1})\Delta x
\]
\[
\text{Right sum: } R_n = \sum_{i=1}^n f(x_i)\Delta x
\]
\[
\text{Midpoint: } M_n = \sum_{i=1}^n f\left(\frac{x_{i-1}+x_i}{2}\right)\Delta x
\]
\[
\Delta x=\frac{b-a}{n},\quad x_i=a+i\Delta x
\]

\section*{Linearization / Tangent Line Approximation}
\[
L(x) = f(a) + f'(a)(x-a)
\]
Approximation:
\[
f(x) \approx L(x)
\]

\section*{Fundamental Theorem of Calculus}
\subsection*{Part 1}
If $F'(x)=f(x)$, then
\[
\frac{d}{dx}\int_a^x f(t)\,dt = f(x).
\]

\subsection*{Part 2}
\[
\int_a^b f(x)\,dx = F(b)-F(a).
\]

\section*{Continuity}
A function is continuous at $x=a$ if:
\[
\lim_{x\to a}f(x)=f(a).
\]

\section*{Closed Interval and Extreme Values}
If $f$ is continuous on $[a,b]$, then it has an absolute max and an absolute min on $[a,b]$.

\section*{Concavity and Inflection Points}
\[
f''(x)>0 \Rightarrow \text{concave up}, \quad f''(x)<0 \Rightarrow \text{concave down}.
\]
Inflection point: where concavity changes and $f''(x)=0$ or undefined.

\section*{Mean Value Theorem}
If $f$ is continuous on $[a,b]$ and differentiable on $(a,b)$, then
\[
f'(c)=\frac{f(b)-f(a)}{b-a}.
\]

\section*{Rolle's Theorem}
If $f(a)=f(b)$ and $f$ satisfies MVT conditions, then 
\[
f'(c)=0
\]
for some $c$ in $(a,b)$.

\section*{Tangent Line Formula}
Slope: $m=f'(a)$  
Equation:  
\[
y-f(a)=f'(a)(x-a)
\]
\section*{Reversing Limits of Integration}
If the bounds of an integral are reversed:
\[
\int_b^a f(x)\,dx = -\int_a^b f(x)\,dx.
\]

This is super helpful when simplifying definite integrals or matching the Fundamental Theorem of Calculus.

\section*{Even and Odd Functions in Integration}

\subsection*{Even Functions}
A function is \textbf{even} if:
\[
f(-x) = f(x).
\]
Examples: $\cos x$, $x^2$, $|x|$

Integral shortcut:
\[
\int_{-a}^{a} f(x)\,dx = 2\int_{0}^{a} f(x)\,dx.
\]

\subsection*{Odd Functions}
A function is \textbf{odd} if:
\[
f(-x) = -f(x).
\]
Examples: $\sin x$, $x^3$, $\tan x$

Integral shortcut:
\[
\int_{-a}^{a} f(x)\,dx = 0.
\]

\section*{Integrals from \(0\) to \(x\) and Variable Upper Limits}

If 
\[
F(x) = \int_0^x f(t)\,dt,
\]
then by the Fundamental Theorem of Calculus:
\[
F'(x) = f(x).
\]

More generally:
\[
\frac{d}{dx}\left( \int_{a}^{g(x)} f(t)\,dt \right)
= f(g(x)) \cdot g'(x).
\]

And if the lower limit varies:
\[
\frac{d}{dx}\left( \int_{g(x)}^{a} f(t)\,dt \right)
= -f(g(x)) \cdot g'(x).
\]

\section*{Integrals From \(-a\) to \(a\)}

Use symmetry when possible:

\[
\int_{-a}^{a} f(x)\,dx =
\begin{cases}
0, & f(x) \text{ odd} \\[6pt]
2\displaystyle\int_0^a f(x)\,dx, & f(x) \text{ even}
\end{cases}
\]

This is especially useful in trig integrals and absolute value integrals.

\section*{Related Rates Templates}

\subsection*{General Steps}
\begin{enumerate}
\item Draw a diagram.
\item Identify all variables.
\item Write an equation relating quantities.
\item Differentiate both sides w.r.t. time $t$.
\item Substitute known values.
\item Solve for the desired rate.
\end{enumerate}

\subsection*{Cone}
Volume:
\[
V=\frac{1}{3}\pi r^2 h.
\]
If radius and height are proportional, use similar triangles before differentiating.

\subsection*{Similar Triangles}
If two triangles are similar:
\[
\frac{x}{a}=\frac{y}{b}.
\]
Differentiate after expressing one variable in terms of the other.

\subsection*{Pythagorean (Ladder)}
\[
x^2+y^2 = L^2.
\]
Differentiate:
\[
2x\frac{dx}{dt}+2y\frac{dy}{dt}=0.
\]

\subsection*{Area Changing}
\[
A = \frac{1}{2}bh,\quad A=\pi r^2,\quad A=s^2.
\]
\section*{Intermediate Value Theorem (IVT)}
If $f$ is continuous on $[a,b]$ and $N$ is between $f(a)$ and $f(b)$, then there exists some $c \in (a,b)$ such that
\[
f(c)=N.
\]

\textbf{Use IVT to show existence only (not to find $c$).}

Steps:
\begin{enumerate}
\item Verify continuity on $[a,b]$
\item Compute $f(a)$ and $f(b)$
\item Show the target value lies between them
\end{enumerate}

\section*{Inverse Trig Integrals}
\begin{align*}
\int \frac{1}{a^2+x^2}\,dx &= \frac{1}{a}\arctan\!\left(\frac{x}{a}\right)+C \\
\int \frac{1}{\sqrt{a^2-x^2}}\,dx &= \arcsin\!\left(\frac{x}{a}\right)+C \\
\int \frac{1}{|x|\sqrt{x^2-a^2}}\,dx &= \arcsec\!\left(\frac{|x|}{a}\right)+C
\end{align*}

\textbf{Tip:} Match the form exactly before integrating.

\section*{Exponential Growth and Decay}
\[
\frac{dP}{dt}=kP
\]

Solution:
\[
P(t)=P_0e^{kt}
\]

\begin{itemize}
\item $k>0$ growth
\item $k<0$ decay
\end{itemize}

Doubling time:
\[
T=\frac{\ln 2}{k}
\]

Half-life:
\[
T=\frac{\ln 2}{|k|}
\]

\section*{Differential Equations}

\subsection*{Separable Equations}
\[
\frac{dy}{dx}=g(x)h(y)
\]

Steps:
\begin{enumerate}
\item Separate variables
\item Integrate both sides
\item Add $C$
\item Solve for $y$
\end{enumerate}

\section*{Initial Value Problems}
Use the given condition (ex: $y(0)=y_0$) to solve for $C$.

\textbf{Important:} Plug in the initial value \emph{after} integrating.
\section*{Trig Substitution}
\begin{center}
\begin{tabular}{c c}
Expression & Substitution \\ \hline
$\sqrt{a^2-x^2}$ & $x=a\sin\theta$ \\
$\sqrt{a^2+x^2}$ & $x=a\tan\theta$ \\
$\sqrt{x^2-a^2}$ & $x=a\sec\theta$
\end{tabular}
\end{center}

Steps:
\begin{enumerate}
\item Substitute
\item Use trig identities
\item Integrate
\item Convert back using a triangle
\end{enumerate}
\section*{Optimization Problems}

Common goals: maximum area, minimum cost, shortest distance.

Steps:
\begin{enumerate}
\item Draw a diagram
\item Write the objective function
\item Use constraints to reduce to one variable
\item Find critical points
\item Check endpoints / reasonableness
\end{enumerate}

Common setups:
\begin{itemize}
\item Fencing problems
\item Box problems
\item Distance problems
\end{itemize}

\end{document}
